%%
%% EPITECH PROJECT, 2025
%% Tek2 Internship Report
%% File description:
%% Mission Outcome chapter
%% of the report
%%

\section{Venues Workflow}

\begin{flushleft}
  Le workflow des sites d'évènements a été implémenté en
  (environ) trois semaines, puis a été déployé et utilisé
  en production chez le client.
\end{flushleft}

\subsection*{Facilités}

\begin{flushleft}
  Le projet était délibérément simple dans son architecture,
  étant donné qu'il s'agissait du premier projet qui avait
  été attribué aux stagiaires. De plus, le projet été à
  réaliser en Python, qui est un langage de programmation
  que je maîtrise déjà assez pour pouvoir écrire du code
  sans devoir réfléchir à ce que je fais.
\end{flushleft}

\subsection*{Difficultés}

\begin{flushleft}
  Malgré les facilités du projet, il y avait évidemment
  aussi des difficultés. La difficulté primaire du projet
  était de pouvoir gérer des grandes tailles de fichiers
  venant de l'utilisateur, car le besoin de mémoire vive,
  ainsi que le temps nécessaire pour faire tout le traitement
  qui vient avec, dépend en grande partie de la taille du fichier.
  Le temps était un grand problème, car il y a une limite dure de
  15 minutes sur les fonctions (il s'agit d'une limitation dans AWS),
  ce qui faisait qu'un traitement excédant cette durée était
  automatiquement annulé.
\end{flushleft}

\begin{flushleft}
  Pour régler ce problème, il fallait diviser le fichier prit en
  entrée en plusieurs petits morceaux (qui eux sont facile à traiter),
  puis de faire le traitement de morceaux en parallèle. Cette approche
  avait nettement réduit la durée du traitement (de \~14 minutes avant
  l'amélioration à \~2 minutes après l'amélioration).
\end{flushleft}

\subsection*{Réception}

\begin{flushleft}
  Peu après le déploiement, JS (le représentant du client) nous a
  confirmé que l'équipe des photographes avait pu utiliser le workflow
  avec succès, et qu'il était bien content de ne plus devoir manipuler
  les fichiers soi-même.
\end{flushleft}

\newpage

\section{MAM Shortcuts Extension}

\begin{flushleft}
  L'extension de navigateur pour les raccourcis de
  clavier a abouti après environ 11 semaines.
  Comme il s'agit le développement lié à un navigateurs,
  j'ai dû réaliser le projet dans le langage de programmation
  JavaScript. Cependant, comme c'est le cas dans l'industrie,
  j'ai utilisé une extension de ce langage nommé \textit{TypeScript},
  qui permet de mieux organiser son code, ainsi que d'éviter des
  fautes de typage. De plus, comme je souhaitais faire une belle interface
  graphique pour les utilisateurs, j'ai choisi de faire du \textit{React},
  qui est aussi une technologie typiquement utilisée dans l'industrie.
\end{flushleft}

\begin{flushleft}
  Quant au status, le projet est terminé, mais il reste améliorable
  avec des nouvelles fonctionnalités, vu qu'il s'agit d'un projet
  ouvert. Lors de ma réalisation, j'ai tâché de ne pas inclure
  d'informations sur le client dans le code, ce qui veut dire qu'il
  serait possible que le client ouvre projet au public un jour ou
  un autre, en fonction de leur politique de confidentialité.
\end{flushleft}

\subsection*{Facilités}

\begin{flushleft}
  Ce qui m'a aidé le plus lors de la réalisation du projet, était
  que pour l'extension de navigateur, j'avais une grande vague de
  motivation et beaucoup d'idées. Comme il s'agissait d'un projet
  relativement ouvert, JS était d'accord avec les idées que j'avais
  quand je les ai présenté lors des réunions hebdomadaires, ce qui
  me motivait d'avantage de travailler sur le projet.
\end{flushleft}

\subsection*{Difficultés}

\begin{flushleft}
  Contrairement au premier projet, je n'avais pas autant d'expérience
  avec le React voire le TypeScript qu'avec le Python, ce qui a ralenti
  le début du projet. De plus, j'ai dû changer du code que j'avais déjà
  écrit, comme c'était un projet relativement ouvert, et que des nouvelles
  idées nécessitant de changer le code étaient assez typique.
  Une autre difficulté était de se limiter à un minimum strict de dépendences
  externes, \sccad{} de code que d'autres développeurs partagent. La raison
  que je me limitait à un minimum sur les dépendences étant d'un côté le
  côté éducatif (il s'agit d'un stage, après tout), mais surtout qu'il y a
  souvent des failles de sécurité venant de dépendances externes, et que je
  voulais éviter cela, connaissant la politique de confidentialité du client.
\end{flushleft}

\subsection*{Réception}

\begin{flushleft}
  Vers la fin du développement du projet, j'avais la chance de pouvoir
  démontrer l'extension de navigateur à toute l'équipe du client lors
  de leur réunion trimestrielle, après avoir été invité par JS.\@
  L'équipe était contente et un peu surprise, car le projet était présenté
  comme bonus à la fin de la partie de développement de la réunion, et que
  qu'ils ne savait pas à quoi s'attendre.
\end{flushleft}

\section{MAM SDK}

\begin{flushleft}
  Le kit de développement de logiciel pour le MAM a été réalisé en Python,
  afin de pouvoir être utilisé dans d'autres projets qui sont crée dans le
  même langage. Le projet était en développement en même temps que le serveur
  MCP, car les outils du serveur MCP sont crée à partir de l'SDK, ce qui veut
  dire qu'on pouvait tout de suite le mettre en pratique et vérifier que tout
  marche comme attendu dans un cadre pratique.
\end{flushleft}

\subsection*{Facilités}

\begin{flushleft}
  Le développement d'outils pour d'autres développeurs est mon domaine
  préféré, donc j'ai eu beaucoup de facilités lors de ce projet.
  De plus, j'ai bien apprécié de lire la spécification et d'en créer
  un outil qui peut être réutilisé dans d'autres projets.
\end{flushleft}

\subsection*{Difficultés}

\begin{flushleft}
  La difficulté principale du projet était de bien nommer les opérations.
  Typiquement, les services numériques disposant d'un API peuvent documenter
  les noms de leurs opérations avec la spécification OpenAPI.\@ Or, le MAM
  du client n'a pas correctement documenté de noms, ce qui faisait que je devais
  générer des noms d'opérations avec les informations que j'avais.
\end{flushleft}

\begin{flushleft}
  Au final, cette solution a bien marchée, et avec un peu d'ajustements via
  un fichier de configuration que j'ai crée, le produit final marche sans
  problèmes.
\end{flushleft}

\subsection*{Réception}

\begin{flushleft}
  Comme il s'agit d'un projet / outil interne aux développeurs, il n'y
  a pas eu d' ``utilisteurs'' du produit à proprement dire, et donc il
  n'y avait pas de vraie ``récéption'' non plus.
  Cependant, comme dit auparavant, on a pu utiliser le projet lors du
  développement du serveur MCP, et tout a fonctionné comme prévu.
\end{flushleft}

\newpage

\section{MAM MCP Server}

\begin{flushleft}
  Le serveur MCP pour interagir avec la plateforme MAM
  était le dernier projet sur lequel nous avions travaillé.
  On a implémenté le projet en Python en utilisant la libraire
  \textit{FastMCP}, qui fait automatiquement toute la partie
  de gestion d'entrée qui est répétitive à implémenter.
  Comme tout ce qui est en rapport avec l'intelligence artificielle
  est assez nouveau, le projet à dû passer par plusieurs itérations
  avant d'avoir un produit final. Ainsi, on devait complètement
  retravailler l'architecture du projet deux reprises, ce qui nous
  a prit du temps.
\end{flushleft}

\begin{flushleft}
  Le projet du serveur MCP est abouti pour l'instant, mais va
  être reprit par les ingénieurs de TrackIt, donc Mathieu BRUJAN
  et Timothé CONIEL, dès l'année 2026.
\end{flushleft}

\subsection*{Facilités}

\begin{flushleft}
  Le début du projet était très rapide, car il y a beaucoup d'outils
  et de logiciels pour travailler avec des intelligences artificielles
  en local, ce qui était idéal pour tester le serveur MCP et ses
  capacités. De plus, quand il y avait des erreurs, c'était assez facile
  de les repérer, car il est possible de voir exactement ce que l'IA pense,
  et donc d'\scou{} l'erreur provient.
\end{flushleft}

\subsection*{Difficultés}

\begin{flushleft}
  \scemaj{}tant donné que le projet utilise de l'intelligence
  artificielle, et que la nature de tout ce qui sort de cette
  IA est assez volatile, j'ai rencontré un grand nombre de messages
  d'erreur sur ce projet. Le problème vient du fait que l'IA peut
  se tromper, et donc provoquer des erreurs, mais contrairement à
  du code traditionnel, on ne peut pas simplement débugger une IA.\@
  Finalement, j'ai apprit avec mes collègues que les erreurs provenaient
  d'un modèle d'intelligence artificelle qui avait été corrigé, mais
  qui n'avait pas été mis\textendash{}à\textendash{}jour par AWS.\@
\end{flushleft}

\begin{flushleft}
  Une autre difficulté était qu'on ne pouvait pas changer de modèle
  d'intelligence artificielle, car le seule meilleure modèle sur la
  plateforme est un modèle Claude, et qu'Anthropic (l'entreprise
  américaine qui a crée Claude) a changé ses conditions d'utilisation
  pour interdire n'importe quelle entreprise étant lié à la Chine d'une
  manière ou d'une autre d'utiliser leurs produits, ce qui nous
  interdissait effectivement de corriger la source d'erreur.
\end{flushleft}

\newpage

\subsection*{Réception}

\begin{flushleft}
  Comme avec l'extension de navigateur, j'ai à nouveau été
  invité à présenter le projet devant toute l'équipe du client
  pour leur dernière réunion trimestrielle de l'année.
\end{flushleft}

\begin{flushleft}
  Malheureusement, dû à la nature volatile du projet, j'ai
  rencontré deux erreurs pendant la démonstration du projet,
  mais ça n'a pas empêché les membres de l'équipe du client
  de voir le potentiel dans le projet, ainsi que de se poser
  des questions sur d'autres cas d'utilisations pour le projet
  (ce qui était le but de JS), donc les problèmes n'ont
  finalement pas trop empiété au succès du serveur MCP.\@
\end{flushleft}
